% -*- coding: utf-8 -*-
% XeLaTeX 编译
\documentclass[a3paper, landscape, 11pt, twocolumn]{article}

% 引入宏包
\usepackage{fontspec}
\usepackage[UTF8]{ctex}
\usepackage{geometry}
\usepackage{listings}
\usepackage{xcolor}
\usepackage{enumitem}
\usepackage{tabularx}
\usepackage{makecell}
\usepackage{fancyhdr}
\usepackage{amssymb}
\usepackage{lastpage}

% 设置页面边距 (A3横向双栏)
\geometry{left=2cm, right=2cm, top=2.5cm, bottom=2.5cm, columnsep=1.8cm}

% 颜色定义
\definecolor{codebg}{RGB}{245,245,245}
\definecolor{keywordcolor}{RGB}{0,0,150}
\definecolor{commentcolor}{RGB}{0,128,0}
\definecolor{stringcolor}{RGB}{163,21,21}

% Java代码块样式设置
\lstdefinestyle{JavaExam}{
    language=Java,
    basicstyle=\ttfamily\small,
    backgroundcolor=\color{codebg},
    keywordstyle=\color{keywordcolor}\bfseries,
    commentstyle=\color{commentcolor},
    stringstyle=\color{stringcolor},
    showstringspaces=false,
    breaklines=true,
    frame=single,
    rulecolor=\color{black!50},
    tabsize=4,
    xleftmargin=1em,
    xrightmargin=1em,
    aboveskip=1em,
    belowskip=1em
}

% 列表样式
\newlist{options}{enumerate}{1}
\setlist[options]{label=\Alph*., nosep, leftmargin=*, itemsep=0.5em}

% 选择题选项命令
% 注意：选择具体使用方式在题目中定义

% 页眉页脚设置
\pagestyle{fancy}
\fancyhf{}
\fancyhead[C]{\Large\bfseries 《Java程序设计》期末考试试卷}
\fancyfoot[C]{第 \thepage\ 页 共 \pageref{LastPage}\ 页}
\fancyfoot[R]{得分：\underline{\hspace{2cm}}}
\fancyfoot[L]{题号：\underline{\hspace{1cm}}}
\renewcommand{\headrulewidth}{0.4pt}

% 填空题横线
\newcommand{\blank}{\underline{\hspace{2cm}}}

\begin{document}

% 试卷头部
\twocolumn[{
			\begin{center}
				\vspace*{-1cm}
				{\huge\bfseries 《Java程序设计》期末考试试卷（C 卷）}\\[0.8em]
				{\large 考试时间：120 分钟 \hspace{2em} 总分：150 分}\\[1em]
				\hrule
				\vspace{0.8em}
				\begin{tabularx}{0.95\linewidth}{|c|X|c|X|c|X|}
					\hline
					\makecell{题号} & \makecell{一                                                                        \\(40分)} & \makecell{二\\(20分)} & \makecell{三\\(20分)} & \makecell{四\\(30分)} & \makecell{五\\(40分)} \\
					\hline
					得分            & \hspace{1.5cm} & \hspace{1.5cm} & \hspace{1.5cm} & \hspace{1.5cm} & \hspace{1.5cm} \\
					\hline
				\end{tabularx}\\[1em]
				{\large 姓名：\underline{\hspace{3cm}} \hspace{1.5em} 学号：\underline{\hspace{4cm}} \hspace{1.5em} 班级：\underline{\hspace{3cm}}}
				\vspace{0.8em}
				\hrule
				\vspace{1em}
			\end{center}
		}]

\section*{一、单项选择题（共 20 小题，每小题 2 分，共 40 分）}
\textit{（在每小题给出的四个选项中，只有一项是符合题目要求的）}

\begin{enumerate}[label=\arabic*., itemsep=1em]
	\item 下列哪项不属于 Java 语言的关键字？
	      \begin{options}
		      \item goto
		      \item const
		      \item null
		      \item native
	      \end{options}

	\item 关于自动装箱和拆箱的描述，正确的是？
	      \begin{options}
		      \item 自动装箱是将基本类型转换为对应的包装类
		      \item 自动拆箱是将基本类型转换为包装类
		      \item 自动装箱发生在基本类型向包装类赋值时，需要显式调用方法
		      \item 自动拆箱不会发生 NullPointerException
	      \end{options}

	\item 下列哪个不是 Java 中合法的数组声明？
	      \begin{options}
		      \item \texttt{int[] a = new int[5];}
		      \item \texttt{int a[] = \{1,2,3\};}
		      \item \texttt{int a[3] = \{1,2,3\};}
		      \item \texttt{int[] a = new int[]{1,2,3};}
	      \end{options}

	\item 在 Java 中，枚举类型的关键字是？
	      \begin{options}
		      \item enum
		      \item Enum
		      \item enumeration
		      \item enumerator
	      \end{options}

	\item 下列关于 Lambda 表达式的说法，错误的是？
	      \begin{options}
		      \item Lambda 表达式可以替代匿名内部类的所有用法
		      \item Lambda 表达式只能用于函数式接口
		      \item Lambda 表达式的参数类型可以省略
		      \item Lambda 表达式可以引用外部变量，但变量必须是 effectively final 的
	      \end{options}

	\item 下列哪个接口是函数式接口（只有一个抽象方法）？
	      \begin{options}
		      \item Runnable
		      \item Comparator
		      \item Serializable
		      \item Cloneable
	      \end{options}

	\item 关于 Stream API 的描述，正确的是？
	      \begin{options}
		      \item Stream 可以重复使用
		      \item Stream 的中间操作会立即执行
		      \item Stream 的终端操作会触发流的执行
		      \item Stream 只能用于集合，不能用于数组
	      \end{options}

	\item 下列哪个类用于表示不可变日期（Java 8 及以上）？
	      \begin{options}
		      \item Date
		      \item Calendar
		      \item LocalDate
		      \item SimpleDateFormat
	      \end{options}

	\item 关于 Optional 类的说法，正确的是？
	      \begin{options}
		      \item Optional 可以包含 null 值
		      \item Optional.of() 方法允许传入 null
		      \item Optional.ofNullable() 允许传入 null
		      \item Optional 主要用于提高性能
	      \end{options}

	\item 在泛型中，\texttt{? extends Number} 表示？
	      \begin{options}
		      \item 类型参数必须是 Number 或 Number 的子类
		      \item 类型参数必须是 Number 或 Number 的父类
		      \item 类型参数必须是 Number
		      \item 类型参数可以是任意类型
	      \end{options}

	\item 关于 Java 注解（Annotation）的描述，正确的是？
	      \begin{options}
		      \item 注解可以影响程序的运行逻辑
		      \item 注解只能用于类和方法
		      \item \texttt{@Override} 注解用于标识重写父类方法
		      \item 自定义注解需要使用 \texttt{@interface} 关键字
	      \end{options}

	\item 下列哪个反射方法可以获取类的所有 public 字段？
	      \begin{options}
		      \item getDeclaredFields()
		      \item getFields()
		      \item getDeclaredField(String name)
		      \item getField(String name)
	      \end{options}

	\item 关于并发包中的 ConcurrentHashMap，正确的是？
	      \begin{options}
		      \item 它是线程安全的，且不允许 null 键和值
		      \item 它的所有方法都使用 synchronized 修饰
		      \item 它继承自 HashMap
		      \item 它不支持高并发
	      \end{options}

	\item 下列关于 volatile 关键字的说法，正确的是？
	      \begin{options}
		      \item volatile 可以保证原子性
		      \item volatile 可以保证可见性
		      \item volatile 可以替代 synchronized
		      \item volatile 修饰的变量不会被线程缓存
	      \end{options}

	\item 在 Java NIO 中，Buffer 的 flip() 方法的作用是？
	      \begin{options}
		      \item 清空缓冲区
		      \item 反转缓冲区，准备读取操作
		      \item 标记当前位置
		      \item 重置位置到标记
	      \end{options}

	\item 关于 try-with-resources 语句，正确的是？
	      \begin{options}
		      \item 可以自动关闭实现了 AutoCloseable 接口的资源
		      \item 只能关闭一个资源
		      \item 不需要 finally 块，但需要 catch 块
		      \item 资源变量必须是 final 的
	      \end{options}

	\item 下列哪个类不是 java.time 包中的？
	      \begin{options}
		      \item LocalDateTime
		      \item ZonedDateTime
		      \item DateTimeFormatter
		      \item DateFormat
	      \end{options}

	\item 关于 Java 正则表达式的说法，正确的是？
	      \begin{options}
		      \item Pattern 类用于匹配字符串
		      \item Matcher 类的 matches() 方法要求整个字符串匹配
		      \item find() 方法只匹配一次
		      \item group() 方法返回整个匹配结果
	      \end{options}

	\item 在 JDBC 中，使用哪个接口调用存储过程？
	      \begin{options}
		      \item Statement
		      \item PreparedStatement
		      \item CallableStatement
		      \item ProcedureStatement
	      \end{options}

	\item 关于 Swing 中 JTable 的说法，正确的是？
	      \begin{options}
		      \item JTable 默认使用 TableModel 存储数据
		      \item JTable 的列宽不能调整
		      \item JTable 不能添加滚动条
		      \item JTable 中的数据不能编辑
	      \end{options}
\end{enumerate}

\section*{二、判断题（共 10 小题，每小题 2 分，共 20 分）}
\textit{（正确的选 A，错误的选 B）}

\begin{enumerate}[label=\arabic*., itemsep=0.8em]
	\item 抽象类可以包含构造方法。 （ \quad A/B \quad ）
	\item Java 中所有类都直接或间接继承自 Object 类。 （ \quad A/B \quad ）
	\item 接口中的 default 方法可以有方法体。 （ \quad A/B \quad ）
	\item 使用 `==` 比较两个字符串时，比较的是字符串的内容。 （ \quad A/B \quad ）
	\item 一个 try 块可以对应多个 catch 块，且父类异常的 catch 块必须放在子类之后。 （ \quad A/B \quad ）
	\item ArrayList 的初始容量是 10。 （ \quad A/B \quad ）
	\item 泛型通配符 `?` 可以用于创建对象，例如 `new ArrayList<>()`。 （ \quad A/B \quad ）
	\item 线程调用 wait() 方法后，会释放持有的对象锁。 （ \quad A/B \quad ）
	\item FileInputStream 和 FileOutputStream 是字节流。 （ \quad A/B \quad ）
	\item 在 UDP 通信中，DatagramPacket 用于封装数据报。 （ \quad A/B \quad ）
\end{enumerate}

\section*{三、填空题（共 5 小题，每空 2 分，共 20 分）}

\begin{enumerate}[label=\arabic*., itemsep=1em]
	\item Java 中，声明包使用 \underline{\hspace{4cm}} 关键字，声明接口使用 \underline{\hspace{4cm}} 关键字。

	\item 面向对象程序设计的基本特征包括封装、继承和 \underline{\hspace{4cm}}。

	\item 异常分为两类：受检查异常（checked exception）和 \underline{\hspace{4cm}} 异常。

	\item 集合框架中，Set 接口的特点是元素 \underline{\hspace{4cm}} 且 \underline{\hspace{4cm}}。

	\item Java 8 中，引入的函数式编程特性包括 Lambda 表达式和 \underline{\hspace{4cm}} API。
\end{enumerate}

\section*{四、程序运行结果题（共 5 小题，每小题 6 分，共 30 分）}
\textit{（阅读下列程序，写出运行后的输出结果）}

\begin{enumerate}[label=\arabic*., itemsep=1em]
	\item
	      \begin{lstlisting}[style=JavaExam]
public class TestA {
    public static void main(String[] args) {
        int x = 5;
        int y = ++x + x++ + x;
        System.out.println(x);
        System.out.println(y);
    }
}
\end{lstlisting}
	      \textbf{输出：}\underline{\hspace{5cm}}

	\item
	      \begin{lstlisting}[style=JavaExam]
class Animal {
    public void sound() {
        System.out.println("Animal sound");
    }
}
class Dog extends Animal {
    public void sound() {
        System.out.println("Bark");
    }
}
public class TestB {
    public static void main(String[] args) {
        Animal a = new Dog();
        a.sound();
    }
}
\end{lstlisting}
	      \textbf{输出：}\underline{\hspace{5cm}}

	\item
	      \begin{lstlisting}[style=JavaExam]
public class TestC {
    public static void main(String[] args) {
        String s = "Java";
        change(s);
        System.out.println(s);
    }
    public static void change(String str) {
        str = "Python";
    }
}
\end{lstlisting}
	      \textbf{输出：}\underline{\hspace{5cm}}

	\item
	      \begin{lstlisting}[style=JavaExam]
import java.util.*;
public class TestD {
    public static void main(String[] args) {
        List<String> list = Arrays.asList("A", "B", "C");
        list.stream()
            .filter(s -> !s.equals("B"))
            .map(s -> s.toLowerCase())
            .forEach(System.out::print);
    }
}
\end{lstlisting}
	      \textbf{输出：}\underline{\hspace{5cm}}

	\item
	      \begin{lstlisting}[style=JavaExam]
public class TestE {
    static int count = 0;
    public static void main(String[] args) {
        try {
            method();
        } catch (StackOverflowError e) {
            System.out.println("Overflow");
        }
        System.out.println(count);
    }
    public static void method() {
        count++;
        method();
    }
}
\end{lstlisting}
	      \textbf{输出：}\underline{\hspace{5cm}}
\end{enumerate}

\newpage

\section*{五、编程题（共 2 小题，每小题 20 分，共 40 分）}
\textit{（要求代码完整、结构清晰、注释合理）}

\begin{enumerate}[label=\arabic*., itemsep=1em]
	\item \textbf{设计一个简单的在线购物系统} \\
	      定义接口 \texttt{Product}，包含方法 \texttt{String getName()}、\texttt{double getPrice()}。\\
	      实现两个类 \texttt{Electronics}（电子产品）和 \texttt{Clothing}（服装），它们都实现 \texttt{Product} 接口。\\
	      \texttt{Electronics} 有额外属性 \texttt{warrantyMonths}（保修月数），\texttt{Clothing} 有额外属性 \texttt{size}（尺寸）。\\
	      定义类 \texttt{ShoppingCart}，内部使用 \texttt{ArrayList<Product>} 存储商品，提供方法：
	      \begin{itemize}
		      \item \texttt{void addProduct(Product p)}
		      \item \texttt{void removeProduct(Product p)}
		      \item \texttt{double calculateTotal()}
		      \item \texttt{void displayItems()}
	      \end{itemize}
	      在测试类 \texttt{Main} 中，创建若干电子产品与服装对象，添加到购物车，计算总价并显示所有商品信息。
	      \vspace{6cm}

	\item \textbf{实现一个图书管理类} \\
	      定义 \texttt{Book} 类，属性包括 ISBN（\texttt{isbn}）、书名（\texttt{title}）、作者（\texttt{author}）、价格（\texttt{price}）。\\
	      提供构造方法、getter/setter 方法，并重写 \texttt{equals()} 和 \texttt{hashCode()} 方法，以 ISBN 为唯一标识。\\
	      定义 \texttt{Library} 类，使用 \texttt{HashSet<Book>} 存储图书集合，提供以下方法：
	      \begin{itemize}
		      \item \texttt{boolean addBook(Book book)}：添加图书，如果 ISBN 重复则添加失败。
		      \item \texttt{boolean removeBook(String isbn)}：根据 ISBN 删除图书。
		      \item \texttt{Book findBookByIsbn(String isbn)}：根据 ISBN 查找图书。
		      \item \texttt{List<Book> findBooksByAuthor(String author)}：根据作者返回图书列表。
		      \item \texttt{void displayAllBooks()}：显示所有图书信息。
	      \end{itemize}
	      在测试类中创建若干 \texttt{Book} 对象，测试上述所有功能。
	      \vspace{6cm}
\end{enumerate}

\begin{center}
	\textbf{—— 试卷结束，祝考试顺利 ——}
\end{center}

\end{document}
